\subsection{Mathematical Proof: Error Bounding via Spectral Separation}
\label{appendix:spectral_separation}

To formally justify the efficacy of our router, we derive the relationship between the spectral separation in Figure~1 and the regret upper bound of the meta-algorithm.

In a contextual bandit, the probability of a routing error (selecting a sub-optimal model) is inversely proportional to the margin of separation in the feature space. Even with only 3.10\% variance in PC1, the bimodal distribution creates a ``low-density'' corridor that acts as a natural buffer for the decision engine.

\subsubsection{The Margin of Decision}

We define the decision boundary $H$ at $\text{PC1} = 0.3$. For any prompt $x$, the distance to this boundary is $d(x, H)$. Routing is stable because the density gap (the ``Cluster Separation'' dashed line) ensures that:
\begin{equation}
P(d(x, H) < \epsilon) \approx 0
\end{equation}
for a sufficiently small $\epsilon$. This means that very few prompts sit in the ``ambiguous zone,'' allowing the Hybrid router to achieve its 99.7\% usage shift without frequent misclassifications.

\subsubsection{Bounding the Routing Regret}

Because the clusters are spatially distinct ($N=1{,}541$ vs. $N=330$), the Thompson Sampling mechanism converges at a rate determined by the sub-Gaussianity of the rewards within each cluster. Given the spectral gap $\Delta_{\text{gap}}$ between the Low and High PC1 clusters:
\begin{equation}
R(T) \le \sum_{a \in \mathcal{A}} \frac{8 \ln(T)}{\Delta_{\text{gap}}} + \left(1 + \frac{\pi^2}{3}\right) \sum_{a \in \mathcal{A}} \Delta_a
\end{equation}
where $\mathcal{A}$ is the set of available models and $\Delta_a$ is the sub-optimality gap for model $a$.

Since $\Delta_{\text{gap}}$ is maximized by the bimodal separation shown in Figure~1, the ``search space'' for the optimal policy is effectively reduced to a binary classification problem rather than a continuous 384-dimensional regression.

\subsubsection{Meta-Algorithm Expected Behavior}

The theoretical Corralling framework motivates our $\eta=1.0$ meta-algorithm design. The standard regret bound for such algorithms takes the form:
\begin{equation}
\label{eq:meta_regret}
R(T) \le \mathcal{O}\left( \frac{\ln(N)}{\eta} + \sqrt{T \ln(T)} \right)
\end{equation}
where $N$ is the number of expert routers in the ensemble, $\eta$ is the learning rate, and $T$ is the time horizon.

This bound suggests sublinear regret growth with two key components:
\begin{itemize}
    \item \textbf{Exploration term} $\frac{\ln(N)}{\eta}$: Captures the cost of identifying the best expert router among $N$ candidates. With $\eta=1.0$, this term is minimized while maintaining stable weight updates.
    \item \textbf{Convergence term} $\sqrt{T \ln(T)}$: Represents the standard Thompson Sampling convergence rate, independent of the learning rate choice.
\end{itemize}

The choice of $\eta=1.0$ balances rapid adaptation (larger $\eta$ reduces the exploration term) with stability (smaller $\eta$ prevents oscillations in expert weights). Our empirical results in Section~5 confirm that $\eta=1.0$ achieves the best performance across all tested configurations, consistent with this theoretical intuition.

\subsubsection{Decision Margin Hypothesis}

The spectral gap $\Delta_{\text{gap}}$ between the bimodal clusters (Figure~1) minimizes the probability of misclassification through the following mechanism:

\paragraph{Spectral Separation and Routing Error.}
The 3.10\% variance captured by PC1 creates a low-density corridor between prompt types, enabling the router to achieve stable routing decisions with minimal ambiguity. Formally, the routing error probability $P_{\text{error}}$ is bounded by:
\begin{equation}
P_{\text{error}} \le \exp\left(-\frac{\Delta_{\text{gap}}^2}{2\sigma^2}\right)
\end{equation}
where $\sigma^2$ is the variance of the reward distribution within each cluster.

The bimodal structure ensures that $\Delta_{\text{gap}}$ is large relative to $\sigma$, making misclassification exponentially unlikely. This spectral separation directly translates to the routing regret bound in Equation~\eqref{eq:meta_regret}, as the effective number of ``hard'' routing decisions is drastically reduced.

\paragraph{Practical Implications.}
The combination of spectral separation and the meta-algorithm's expected behavior provides both theoretical motivation and empirical validation for our approach. The router can reliably distinguish between prompts requiring different model capabilities, while the meta-algorithm efficiently learns to weight expert opinions without excessive exploration cost.

